\paragraph{Counting the short and boundary sections.}
The orders on the slice control its finite contribution. Its archimedean rank estimate now requires a count of the sections that remain short and of those lost to the boundary. On a connected unfolded complex component put
\(S_i=L_i\) for the \(s\) cutting sections and write
\(w_0=L_0\) for the nonvanishing reference at their
common zeros. The flat-section hypotheses are that these
zeros consist, modulo the uniformizing group, of
\(\delta_k\) distinct framed points, that
\(u_i=S_i/w_0\) give analytic local coordinates there, and
that the chosen short weight provides interpolation of the
values at those points. The interpolation with the required
vanishing of low Fourier coefficients will be used also for
the lower bound on the selected subspace. Both interpolation conditions use the same point count \(\delta=\sum_k\delta_k\) as the heights.

\begin{proposition}[Ranks of the full and short evaluation maps on the slice]
\label{prop:slice-rank-shares}
Assume the simultaneous cut consists on component \(k\) of
\(\delta_k\) distinct framed points at which \(w_0\) is nonzero
and the \(S_i/w_0\) are local coordinates. Assume interpolation
of values at all these points in the short weight with the
prescribed low modes zero. Put \(\delta=\sum_k\delta_k\). For the full rank \(R_j\), the rank deficit \(r_q\) of the rational good space, and the rank \(r_z\) of evaluation on the boundary,
\[
 R_j\le\delta\binom{T+2s+6}{s},\qquad \max(r_q,r_z)\le\frac{22s^2}{T}R_j\quad(s\ge2,\ T\ge11s^2).
\]
 The definite final lift
uses the full space as its good space.
\end{proposition}

\begin{proof}
\proofstep{An upper bound for the entire rank.}
To bound the full rank from above put \(e=T+s+6\).
The jets of order \(\le e\) at the \(\delta_k\) points have
dimension \(\delta_k\binom{e+s}{s}\). A nonzero section
\(v\) in the kernel vanishes to order \(e+1\) there.
We give the maximum argument which excludes it. In the
indefinite chart, use \(b_2,b_1,\xi\) of respective weights
\((s+4)w_L-w_F,\ 2w_L-w_d,\ n w_d\), where
\(n=i g_0^*/2>0\) in this analytic normalization.
All three are nonzero on the component. Put \(S=\sum_i
\|S_i\|^2\), the sum over the \(s\) cutting sections,
and consider
\[
 \Phi=\frac{\|v\|^2\|b_2\|^2\|b_1\|^2
                     \|\xi\|^{2/n}}{S^{e+1/2}} .
\]

The bounded section and the auxiliary factors satisfy the
chart estimates for the first minimum. Near a cut point \(S\gg|u|^2\) and \(v=O(|u|^{e+1})\);
the other numerator factors are bounded there. Thus
\[
 \Phi\ll \frac{|u|^{2(e+1)}}{|u|^{2(e+1/2)}}=|u|.
\] Away from neighborhoods of the cut
orbits the chart estimate gives \(S\gg e^{-2cP}\), and the
numerator is \(O(e^{-4P/g_0^*})\). The hypotheses \(s+6+1/2<T\) and \(Tcg_0^*\le1\) give the strictly positive net decay exponent
\[
  \frac{4}{g_0^*}-2c(e+1/2)>0.
\]
 As the first minimum
goes to zero, \(P\) tends to infinity. Thus every positive
superlevel has compact representatives, away from the cut
points.

On a definite component compactness itself supplies
this step with the bounded auxiliary factors in its chart.
A common nonvanishing locus for the numerator and at
least one cutting section exists by the identity theorem
on the connected domain. In the case \(\delta_k=0\) the
denominator has no common zero and the same argument applies.

A positive maximum contradicts the Levi form. In a local
holomorphic frame the weight of the numerator terms totals
\(e w_L\), leaving a positive half of the strictly positive
light-weight potential in \(\log\Phi\). The log absolute
values of its individual holomorphic factors are
pluriharmonic at a positive maximum.
For a holomorphic factor \(h\) nonzero at the maximum this vanishing reads
\[
  {\cal L}(\log|h|)=0\qquad\text{locally on }\{h\ne0\}.
\]
 The remaining log
sum of squares contributes the pullback of a
Fubini--Study form by the map \([S_1:\cdots:S_s]\).
This maps an \(s\)-dimensional tangent space to one of
dimension \(s-1\).

Choose a nonzero vector in the kernel
of its differential. The Levi form of the log sum is zero
on this vector and that of the remaining potential is
strictly positive. At a maximum their sum must be
nonpositive. This proves injectivity of the jet readings.
Restriction of the global space to all components is
injective, hence summing gives
\[
                      R_j\le\delta\binom{T+2s+6}{s}.
\]

\proofstep{The rank of the complementary sections.}
For the complementary rank take forms \(v_j\) of weight
\(w_F\) interpolating the point delta functions on one
component, with their required low Fourier coefficients
zero. For \(|\alpha|\le T\) the products
\(v_j w_0^{T-|\alpha|}\prod_i S_i^{\alpha_i}\)
belong to the prescribed analytic subspace. This uses
stability of the low-mode vanishing condition under
multiplication by the bounded light forms, whose Fourier
indices are nonnegative. Independence follows by first
taking the least total degree appearing in a relation,
then taking that homogeneous Taylor coefficient at each
point after dividing by \(w_0^T\). Its monomials are
distinct and the value matrix of the \(v_j\) is invertible.
One can extend the families separately by zero to the
other components. The \hyperref[par:shifted-good]{descent of the low-mode space} gives

\[
  \dim H_g\ge\delta\binom{T+s}{s}.
\]
 This identifies
\(r_q=R_j-\dim H_g\) after rationally spanning the cusp
lattice. At the definite final lift the comparison uses
all sections.

\proofstep{Short restrictions and the boundary rank.}
Likewise the \(|\alpha|<T-s\) products from the short
restrictions are independent modulo the span of all
boundary products. Apply the least-degree argument to
a relation of degree less than \(T-s\) against sections
vanishing to that order. There are \(R\delta\) of these
independent classes. The boundary products can be tested
by their rank on the generic coefficient functions: after
extension of scalars their minors vanish precisely when
the corresponding generic complex minors vanish. Taking
a basis out of their short generating family gives

\[
  R\delta+r_z\le R_j.
\]
 This calculation also deals with a
zero-dimensional vector space of readings.

For \(s\ge2\) and \(T\ge11s^2\), put \(B=\binom{T+2s+6}{s}\).
The fraction \(R/B\) is the product of the \(s\) factors
\(1-(2s+7)/(T+2s+6-i)\) for \(0\le i<s\).
A product of \(1-x_i\), \(0\le x_i\le1\), is at least
\(1-\sum x_i\). Consequently
\[
 0\le 1-R/B\le s(2s+7)/T\le22s^2/T .
\]
Since \(R_j/B\le\delta\), both deficient ranks are at most
\(22s^2R_j/T\). This implies the weaker pre-final bound
\(r_z\le r_q+22s^2R_j/T\) used in the slope sum. It uses
no assumption that different primes admit independent
choices of the short family. \label{cor:slice-shares-1}
\end{proof}
